\ifdefined\gkver\else\def\gkver{student}\fi
\documentclass[\gkver]{gaokaozhenti}
\newcommand{\ccnum}[1]{{\fontspec{Noto Sans CJK JP}#1}}
\begin{document}

\chapter{2017年全国理综Ⅲ}
%% paperType: 理综物理部分

\begin{enumerate}
\item
%% number: 14
%% paperName: 2017年全国理综Ⅲ第 14 题 6 分
%% typeId: 1
%% type: 单选题
%% chapter: 曲线运动
%% point: 人造地球卫星
%% method: 无
%% score: 6
%% degree: 850
%% duplicateId: 0
%% body:
2017 年 4 月，我国成功发射的天舟一号货运飞船与天宫二号空间实验室完成了首次交会对接，对接形成的组合体仍沿天宫二号原来的轨道（可视为圆轨道）运行。与天宫二号单独运行相比，组合体运行的\xzanswer{C}
\fourchoices[answer=C]
{周期变大}
{速率变大}
{动能变大}
{向心加速度变大}

%% answer: C
%% memo:
\memoanswer{
由于“天舟一号”与“天宫二号”空间实验室对接形成的组合体仍沿“天宫二号”原来的轨道运行，故轨道半径不变，由万有引力提供向心力得 $G\frac{Mm}{R^{2}}=ma=m\frac{v^{2}}{R}=m\frac{4\pi^{2}}{T^{2}}R$，得 $T=2\pi\sqrt{\frac{R^{3}}{GM}}$，$a=\frac{GM}{R^{2}}$，$v=\sqrt{\frac{GM}{R}}$，故 A、B、D 错误；由于对接形成组合体后质量增加，所以动能增大，故 C 正确。
}

\item
%% number: 15
%% paperName: 2017年全国理综Ⅲ第 15 题 6 分
%% typeId: 1
%% type: 单选题
%% chapter: 电磁感应
%% point: 二级电磁感应
%% method: 无
%% score: 6
%% degree: 550
%% duplicateId: 0
%% body:
如图，在方向垂直于纸面向里的匀强磁场中有一 U 形金属导轨，导轨平面与磁场垂直。金属杆 PQ 置于导轨上并与导轨形成闭合回路 PQRS，一圆环形金属框 T 位于回路围成的区域内，线框与导轨共面。现让金属杆 PQ 突然向右运动，在运动开始的瞬间，关于感应电流的方向，下列说法正确的是\xzanswer{D}
\onepicture[w=4.2cm]{figs/15.png}
\fourchoices[answer=D]
{PQRS 中沿顺时针方向，T 中沿逆时针方向}
{PQRS 中沿顺时针方向，T 中沿顺时针方向}
{PQRS 中沿逆时针方向，T 中沿逆时针方向}
{PQRS 中沿逆时针方向，T 中沿顺时针方向}

%% answer: D
%% memo:
\memoanswer{
由于金属杆 PQ 突然向右运动，导致金属导轨与金属杆 PQ 所围成的面积增大，磁通量增大，由楞次定律知，感应电流产生的磁场阻碍原磁场的变化，故感应电流产生的磁场方向应垂直于纸面向外，PQRS 中的感应电流沿逆时针方向。对于圆环形金属线框 T，金属杆由于运动产生的感应电流所产生的磁场使得 T 内的磁场的磁感应强度变小，磁通量减小，故线框 T 中感应电流产生的磁场方向垂直于纸面向里，故 T 中的感应电流沿顺时针方向，故 D 正确。
}

\item
%% number: 16
%% paperName: 2017年全国理综Ⅲ第 16 题 6 分
%% typeId: 1
%% type: 单选题
%% chapter: 功和能
%% point: 功能原理
%% method: 无
%% score: 6
%% degree: 650
%% duplicateId: 0
%% body:
如图，一质量为 $m$，长度为 $l$ 的均匀柔软细绳 PQ 竖直悬挂。用外力将绳的下端 Q 缓慢地竖直向上拉起至 M 点，M 点与绳的上端 P 相距 $\frac{1}{3}l$。重力加速度大小为 $g$。在此过程中，外力做的功为\xzanswer{A}
\onepicture[w=2.8cm]{\includesvg{figs/16}}
\fourchoices[answer=A]
{$\frac{1}{9}mgl$}
{$\frac{1}{6}mgl$}
{$\frac{1}{3}mgl$}
{$\frac{1}{2}mgl$}

%% answer: A
%% memo:
\memoanswer{
由题意知，将 Q 点拉至 M 点时，外力对细绳所做的功等于细绳 MQ 段克服重力所做的功，将 Q 点拉至 M 点的过程中，细绳 MQ 段重心上升的高度 $\Delta h=\left(\frac{l}{3}+\frac{2l}{3}\times\frac{1}{2}\right)-\left(\frac{l}{3}+\frac{2l}{3}\times\frac{1}{2}\times\frac{1}{2}\right)=\frac{1}{6}l$（取悬挂点为基准），故外力做的功 $W=\frac{2}{3}mg\Delta h=\frac{1}{9}mgl$，故 A 正确。
}

\item
%% number: 17
%% paperName: 2017年全国理综Ⅲ第 17 题 6 分
%% typeId: 1
%% type: 单选题
%% chapter: 力物体的平衡
%% point: 弹力
%% method: 无
%% score: 6
%% degree: 550
%% duplicateId: 0
%% body:
一根轻质弹性绳的两端分别固定在水平天花板上相距 $80\Ucm$ 的两点上，弹性绳的原长也为 $80\Ucm$。将一钩码挂在弹性绳的中点，平衡时弹性绳的总长度为 $100\Ucm$；再将弹性绳的两端缓慢移至天花板上的同一点，则弹性绳的总长度变为（弹性绳的伸长始终处于弹性限度内）\xzanswer{B}
\fourchoices[answer=B]
{$86\Ucm$}
{$92\Ucm$}
{$98\Ucm$}
{$104\Ucm$}

%% answer: B
%% memo:
\memoanswer{
将一钩码挂在弹性绳的中点，平衡后对钩码进行受力分析，如图所示，设钩码质量为 $m$，弹性绳的拉力为 $T$，弹性绳与竖直方向的夹角为 $\theta$，由平衡条件知 $2T\cos\theta=mg$，由胡克定律得 $T=k\Delta x$，联立解得 $k=\frac{mg}{2\cos\theta\Delta x}$，将弹性绳两端拉至同一点重新平衡后，设弹性绳上拉力为 $T_{1}$，弹性绳的形变量为 $\Delta x'$，由平衡条件得 $2T_{1}=mg$，由胡克定律得 $T_{1}=k\Delta x'$，联立解得 $k=\frac{mg}{2\Delta x}$。由几何关系得 $\cos\theta=\frac{3}{5}$，$\Delta x=0.2\Um$，联立解得 $\Delta x'=0.12\Um$，故弹性绳的总长度变为 $92\Ucm$，故 B 正确。

\onepicture[w=6.0cm]{figs/17_an.png}
}

\item
%% number: 18
%% paperName: 2017年全国理综Ⅲ第 18 题 6 分
%% typeId: 1
%% type: 单选题
%% chapter: 磁场
%% point: 磁场的叠加
%% method: 无
%% score: 6
%% degree: 550
%% duplicateId: 0
%% body:
如图，在磁感应强度大小为 $B_{0}$ 的匀强磁场中，两长直导线 P 和 Q 垂直于纸面固定放置，两者之间的距离为 $l$。在两导线中均通有方向垂直于纸面向里的电流 $I$ 时，纸面内与两导线距离为 $l$ 的 a 点处的磁感应强度为零。如果让 P 中的电流反向、其他条件不变，则 a 点处磁感应强度的大小为\xzanswer{C}
\onepicture[w=3.8cm]{figs/18.png}
\fourchoices[answer=C]
{0}
{$\frac{\sqrt{3}}{3}B_{0}$}
{$\frac{2\sqrt{3}}{3}B_{0}$}
{$2B_{0}$}

%% answer: C
%% memo:
\memoanswer{
设导线 P 在点 a 处产生的磁感应强度为 $B$，由于 a 处的磁感应强度为零，故导线 P、Q 在 a 点处产生的合磁感应强度与 $B_{0}$ 等大反向。如图甲所示，由几何关系得，导线 P、Q 在 a 点处产生的合磁感应强度 $B_{0}=2B\cos30^{\circ}$，方向水平向右。若 P 中电流反向、其他条件不变，如图乙所示，由几何关系得，P、Q 导线在 a 点处磁感应强度变为 $B$，方向竖直向上，则 a 点磁感应强度的大小为 $\sqrt{B_{0}^{2}+\left(\frac{B_{0}}{\sqrt{3}}\right)^{2}}=\frac{2\sqrt{3}}{3}B_{0}$，故 C 正确。

\twopicture[numA=甲, numB=乙, labelprefix={}, wA=5.5cm, wB=5.5cm]{figs/18_an1.png}{figs/18_an2.png}
}

\item
%% number: 19
%% paperName: 2017年全国理综Ⅲ第 19 题 6 分
%% typeId: 2
%% type: 多选题
%% chapter: 光学
%% point: 光电效应
%% method: 无
%% score: 6
%% degree: 550
%% duplicateId: 0
%% body:
在光电效应实验中，分别用频率为 $\nu_{a}$、$\nu_{b}$ 的单色光 a、b 照射到同种金属上，测得相应的遏止电压分别为 $U_{a}$ 和 $U_{b}$、光电子的最大初动能分别为 $E_{ka}$ 和 $E_{kb}$。$h$ 为普朗克常量。下列说法正确的是\xzanswer{BC}
\fourchoices[answer=BC]
{若 $\nu_{a}>\nu_{b}$，则一定有 $U_{a}<U_{b}$}
{若 $\nu_{a}>\nu_{b}$，则一定有 $E_{ka}>E_{kb}$}
{若 $U_{a}<U_{b}$，则一定有 $E_{ka}<E_{kb}$}
{若 $\nu_{a}>\nu_{b}$，则一定有 $h\nu_{a}-E_{ka}>h\nu_{b}-E_{kb}$}

%% answer: BC
%% memo:
\memoanswer{
由光电效应方程 $E_{k}=h\nu-W_{0}$ 可知，若 $\nu_{a}>\nu_{b}$，则 $E_{ka}>E_{kb}$，遏止电压用于反映最大初动能的大小，$eU=E_{k}$，则 $U_{a}>U_{b}$，故 A 错误、B 正确；若 $U_{a}<U_{b}$，则 $E_{ka}<E_{kb}$，故 C 正确；$W_{0}=h\nu-E_{k}$，金属的逸出功由金属材料本身决定，故 D 错误。
}

\item
%% number: 20
%% paperName: 2017年全国理综Ⅲ第 20 题 6 分
%% typeId: 2
%% type: 多选题
%% chapter: 运动和力
%% point: 牛顿定律中的图象
%% method: 无
%% score: 6
%% degree: 550
%% duplicateId: 0
%% body:
一质量为 $2\Ukg$ 的物块在合外力 $F$ 的作用下从静止开始沿直线运动。$F$ 随时间 $t$ 变化的图线如图所示，则\xzanswer{AB}
\onepicture[w=6.5cm]{\includesvg{figs/20}}
\fourchoices[answer=AB]
{$t=1\Us$ 时物块的速率为 $1\Ums$}
{$t=2\Us$ 时物块的动量大小为 $4\Ukgms$}
{$t=3\Us$ 时物块的动量大小为 $5\Ukgms$}
{$t=4\Us$ 时物块的速度为零}

%% answer: AB
%% memo:
\memoanswer{
由动量定理可知 $I=Ft=mv$，则 $t=1\Us$ 时物块的速度大小为 $v_{1}=1\Ums$，故 A 正确；$t=2\Us$ 时物块的动量大小 $p_{2}=I_{2}=4\Ukgms$，故 B 正确；$t=3\Us$ 时物块的动量大小 $p_{3}=I_{3}=2\times2\Ukgms-1\times1\Ukgms=3\Ukgms$，故 C 错误；$t=4\Us$ 时物块的速度大小 $v_{4}=\frac{I_{4}}{m}=\frac{2\times2-1\times2}{2}\Ums=1\Ums$，故 D 错误。
}

\item
%% number: 21
%% paperName: 2017年全国理综Ⅲ第 21 题 6 分
%% typeId: 2
%% type: 多选题
%% chapter: 电场
%% point: 等势面
%% method: 无
%% score: 6
%% degree: 450
%% duplicateId: 0
%% body:
一匀强电场的方向平行于 $xOy$ 平面，平面内 a、b、c 三点的位置如图所示，三点的电势分别为 $10\UV$、$17\UV$、$26\UV$。下列说法正确的是\xzanswer{ABD}
\onepicture[w=5.8cm]{\includesvg{figs/21}}
\fourchoices[answer=ABD]
{电场强度的大小为 $2.5\UVcm$}
{坐标原点处的电势为 $1\UV$}
{电子在 a 点的电势能比在 b 点的低 $7\UeV$}
{电子从 b 点运动到 c 点，电场力做功为 $9\UeV$}

%% answer: ABD
%% memo:
\memoanswer{
在匀强电场中，两平行且等长的线段两端的电势差相等，设 O 点电势为 $\varphi$，则 $U_{aO}=U_{cb}$，故 $\varphi=1\UV$，故 B 正确；连接 O、c，过点 a 作 Oc 的垂线，分析可知此垂线为一等势线，则 Oc 沿电场线方向，$E=\frac{26-1}{\sqrt{6^{2}+8^{2}}}\UVcm=2.5\UVcm$，故 A 正确。

\onepicture[w=7.0cm]{figs/21_an.png}

电子在 a 点的电势能 $E_{pa}=-10\UeV$，在 b 点电势能 $E_{pb}=-17\UeV$，$E_{pa}-E_{pb}=7\UeV$，电子在 a 点的电势能比 b 点的高 $7\UeV$，故 C 错误；电子在 c 处的电势能 $E_{pc}=-26\UeV$，电子从 b 点运动到 c 点，电场力的功 $W=-\Delta E_{p}=-(E_{pc}-E_{pb})=9\UeV$，故 D 正确。
}

\item
%% number: 22
%% paperName: 2017年全国理综Ⅲ第 22 题 6 分
%% typeId: 4
%% type: 实验题
%% chapter: 力物体的平衡
%% point: 验证平行四边形实验
%% method: 无
%% score: 6
%% degree: 650
%% duplicateId: 0
%% body:
某探究小组做“验证力的平行四边形定则”实验，将画有坐标轴（横轴为 $x$ 轴，纵轴为 $y$ 轴，最小刻度表示 $1\Umm$）的纸贴在桌面上，如图（a）所示。将橡皮筋的一端 Q 固定在 $y$ 轴上的 B 点（位于图示部分除外），另一端 P 位于 $y$ 轴上的 A 点时，橡皮筋处于原长。
\begin{enumerate}
	\item
	用一只测力计将橡皮筋的 P 端沿 $y$ 轴从 A 点拉至坐标原点 $O$，此时拉力 $F$ 的大小可由测力计读出。测力计的示数如图（b）所示，$F$ 的大小为\tkanswer[1cm]{$4.0\UN$}。
	\item
	撤去（1）中的拉力，橡皮筋 P 端回到 A 点；现使用两个测力计同时拉橡皮筋，再次将 P 端拉至 O 点，此时观察到两个拉力分别沿图（a）中两条虚线所示的方向，由测力计的示数读出两个拉力的大小分别为 $F_{1}=4.2\UN$ 和 $F_{2}=5.6\UN$。
	\begin{enumerate}
		\item
		用 $5\Umm$ 长度的线段表示 $1\UN$ 的力，以 O 点为作用点，在图（a）中画出力 $F_{1}$、$F_{2}$ 的图示，然后按平行四边形定则画出它们的合力；
		\item
		合力的大小为\tkanswer[1cm]{$4.0\UN$}，合力与拉力 $F$ 的夹角的正切值为\tkanswer[1cm]{0.05}。若合力与拉力 $F$ 的大小及方向的偏差均在实验所允许的误差范围之内，则该实验验证了力的平行四边形定则。
	\end{enumerate}
\end{enumerate}
\twopicture[numA={（a）}, numB={（b）}, labelA=2017全国理综Ⅲ22a, labelB=2017全国理综Ⅲ22b, align=right, wA=7.2cm, wB=5.6cm]
{figs/22a.png}
{figs/22b.png}

%% answer: 
\jdanswer{
\begin{enumerate}
	\item
	4.0

	\item
	（i）如图所示。（ii）4.0，0.05

\end{enumerate}
}
%% memo:
\memoanswer{
（1）由题图可知，$F$ 的大小为 $4.0\UN$。

（2）（i）由题意画出 $F_{1}$、$F_{2}$ 的图示，如图所示，$F_{1}$ 用长为 $21\Umm$ 的线段表示，$F_{2}$ 用长为 $28\Umm$ 的线段表示。

\onepicture[w=6.0cm]{figs/22_an.png}

（ii）由图示，测得合力 $F_{\text{合}}$ 的长度为 $20\Umm$，则 $F_{\text{合}}$ 的大小为 $4.0\UN$，利用作图法可得，$F_{\text{合}}$ 与 $F$ 的正切值为 0.05。
}

\item
%% number: 23
%% paperName: 2017年全国理综Ⅲ第 23 题 9 分
%% typeId: 4
%% type: 实验题
%% chapter: 电路
%% point: 多用表的使用
%% method: 无
%% score: 9
%% degree: 350
%% duplicateId: 0
%% body:
图（a）为某同学组装完成的简易多用电表的电路图。图中 $E$ 是电池；$R_{1}$、$R_{2}$、$R_{3}$、$R_{4}$ 和 $R_{5}$ 是固定电阻，$R_{6}$ 是可变电阻；表头 G 的满偏电流为 $250\UuA$，内阻为 $480\UO$。虚线方框内为换挡开关，A 端和 B 端分别与两表笔相连。该多用电表有 5 个挡位，5 个挡位为：直流电压 $1\UV$ 挡和 $5\UV$ 挡，直流电流 $1\UmA$ 挡和 $2.5\UmA$ 挡，欧姆$\times 100\UO$ 挡。
\begin{enumerate}
	\item
	图（a）中的 A 端与\tkanswer[1cm]{黑}（填“红”或“黑”）色表笔相连接。
	\item
	关于 $R_{6}$ 的使用，下列说法正确的是\tkanswer[1cm]{B}（填正确答案标号）。
	\threechoices[answer=B]
	{在使用多用电表之前，调整 $R_{6}$ 使电表指针指在表盘左端电流“0”位置}
	{使用欧姆挡时，先将两表笔短接，调整 $R_{6}$ 使电表指针指在表盘右端电阻“0”位置}
	{使用电流挡时，调整 $R_{6}$ 使电表指针尽可能指在表盘右端电流最大位置}
	\item
	根据题给条件可得 $R_{1}+R_{2}=$\tkanswer[1.5cm]{$160\UO$}，$R_{4}=$\tkanswer[1.5cm]{$880\UO$}。
	\item
	某次测量时该多用电表指针位置如图（b）所示。若此时 B 端是与“1”连接的，则多用电表读数为\tkanswer[2cm]{$1.47\UmA$}；若此时 B 端是与“3”连接的，则读数为\tkanswer[2.5cm]{$1.10\times10^{3}\UO$}；若此时 B 端是与“5”连接的，则读数为\tkanswer[2cm]{$2.95\UV$}。（结果均保留 3 位有效数字）
\end{enumerate}
\twopicture[numA={（a）}, numB={（b）}, labelA=2017全国理综Ⅲ23a, labelB=2017全国理综Ⅲ23b, align=right, wA=6.2cm, wB=8.0cm]
{figs/23a.png}
{figs/23b.png}

%% answer: 
\jdanswer{
\begin{enumerate}
	\item
	黑

	\item
	B

	\item
	160；880

	\item
	$1.47\UmA$；$1.10\times10^{3}\UO$；$2.95\UV$

\end{enumerate}
}
%% memo:
\memoanswer{
（1）由“红进黑出”的原则，可知 A 端与黑表笔连接。

（2）多用电表只有接欧姆挡时才会用到内部的电源，故多用电表接 3 时，通过调节滑动变阻器达到欧姆调零的目的，故 B 正确。

（3）当多用电表接“2”时为小量程电流表，量程为 $1\UmA$，表头满偏时为 $250\UuA$，则经过 $R_{1}$ 和 $R_{2}$ 的电流为 $750\UuA$，由部分电路欧姆定律可得 $R_{1}+R_{2}=\frac{480\UO\times250\UuA}{750\UuA}=160\UO$；当多用电表接“4”时，为小量程电压表，量程为 $1\UV$，则电表此时的总电阻为 $R=\frac{1\UV}{1\times10^{-3}\UA}=1000\UO$，而表头和电阻 $R_{1}$ 和 $R_{2}$ 并联的电阻为 $R'=\frac{480\times160}{480+160}\UO=120\UO$，故 $R_{4}=R-R'=880\UO$。

（4）若 B 端与“1”相连，多用电表为量程为 $2.5\UmA$ 的电流表，读数为 $1.47\UmA$；若 B 端与“3”相连，多用电表为 $\times$100 欧姆挡，则读数为 $1100\UO$，记为 $1.10\times10^{3}\UO$；若 B 端与“5”相连，多用电表为量程为 $5\UV$ 的电压表，则读数为 $2.95\UV$。
}

\item
%% number: 24
%% paperName: 2017年全国理综Ⅲ第 24 题 12 分
%% typeId: 6
%% type: 计算题
%% chapter: 磁场
%% point: 带电粒子在磁场中的运动
%% method: 无
%% score: 12
%% degree: 450
%% duplicateId: 0
%% body:
如图，空间存在方向垂直于纸面（$xOy$ 平面）向里的磁场。在 $x\geq0$ 区域，磁感应强度的大小为 $B_{0}$；$x<0$ 区域，磁感应强度的大小为 $\lambda B_{0}$（常数 $\lambda>1$）。一质量为 $m$、电荷量为 $q$（$q>0$）的带电粒子以速度 $v_{0}$ 从坐标原点 $O$ 沿 $x$ 轴正向射入磁场，此时开始计时，当粒子的速度方向再次沿 $x$ 轴正向时，求（不计重力）：
\begin{enumerate}
	\item
	粒子运动的时间；
	\item
	粒子与 $O$ 点间的距离。
\end{enumerate}
\onepicture[w=6.0cm, align=right]{figs/24.png}

%% answer: 
\jdanswer{
\begin{enumerate}
	\item
	$\frac{\pi m}{qB_{0}}\left(1+\frac{1}{\lambda}\right)$

	\item
	$d=\frac{2mv_{0}}{qB_{0}}\left(1-\frac{1}{\lambda}\right)$

\end{enumerate}
}
%% memo:
\memoanswer{
粒子的运动轨迹如图所示。

\onepicture[w=4.5cm]{figs/24_an.png}

（1）在匀强磁场中，带电粒子做匀速圆周运动。设在 $x\geq0$ 区域，圆周半径 $R_{1}$；在 $x<0$ 区域，圆周半径为 $R_{2}$。由洛伦兹力公式及牛顿第二定律得

$qB_{0}v_{0}=m\frac{v_{0}^{2}}{R_{1}}$\ccnum{①}，$q\lambda B_{0}v_{0}=m\frac{v_{0}^{2}}{R_{2}}$\ccnum{②}，

粒子速度方向转过 $180^{\circ}$ 时，所需时间 $t_{1}$ 为

$t_{1}=\frac{\pi R_{1}}{v_{0}}$\ccnum{③}，

粒子再转过 $180^{\circ}$ 时，所需时间 $t_{2}$ 为

$t_{2}=\frac{\pi R_{2}}{v_{0}}$\ccnum{④}，

联立\ccnum{①}\ccnum{②}\ccnum{③}\ccnum{④}式得，所求时间为

$t_{0}=t_{1}+t_{2}=\frac{\pi m}{qB_{0}}\left(1+\frac{1}{\lambda}\right)$\ccnum{⑤}。

（2）由几何关系及\ccnum{①}\ccnum{②}式得，所求距离为

$d_{0}=2(R_{1}-R_{2})=\frac{2mv_{0}}{qB_{0}}\left(1-\frac{1}{\lambda}\right)$\ccnum{⑥}。
}

\item
%% number: 25
%% paperName: 2017年全国理综Ⅲ第 25 题 20 分
%% typeId: 6
%% type: 计算题
%% chapter: 运动和力
%% point: 牛顿定律的应用
%% method: 无
%% score: 20
%% degree: 350
%% duplicateId: 0
%% body:
如图，两个滑块 A 和 B 的质量分别为 $m_{A}=1\Ukg$ 和 $m_{B}=5\Ukg$，放在静止于水平地面上的木板的两端，两者与木板间的动摩擦因数均为 $\mu_{1}=0.5$；木板的质量为 $m=4\Ukg$，与地面间的动摩擦因数为 $\mu_{2}=0.1$。某时刻 A、B 两滑块开始相向滑动，初速度大小均为 $v_{0}=3\Ums$。A、B 相遇时，A 与木板恰好相对静止。设最大静摩擦力等于滑动摩擦力，取重力加速度大小 $g=10\Umsq$。求：
\begin{enumerate}
	\item
	B 与木板相对静止时，木板的速度；
	\item
	A、B 开始运动时，两者之间的距离。
\end{enumerate}
\onepicture[w=7.0cm, align=right]{\includesvg{figs/25}}

%% answer: 
\jdanswer{
\begin{enumerate}
	\item
	$1\Ums$

	\item
	$1.9\Um$

\end{enumerate}
}
%% memo:
\memoanswer{
（1）滑块 A 和 B 在木板上滑动时，木板也在地面上滑动。设 A、B 和木板所受的摩擦力大小分别为 $f_{1}$、$f_{2}$、$f_{3}$，A 和 B 相对于地面的加速度大小分别为 $a_{A}$ 和 $a_{B}$，木板相对于地面的加速度大小为 $a_{1}$。在滑块 B 与木板达到共同速度前有

$f_{1}=\mu_{1}m_{A}g$\ccnum{①}，$f_{2}=\mu_{1}m_{B}g$\ccnum{②}，$f_{3}=\mu_{2}(m+m_{A}+m_{B})g$\ccnum{③}，

由牛顿第二定律得

$f_{1}=m_{A}a_{A}$\ccnum{④}，$f_{2}=m_{B}a_{B}$\ccnum{⑤}，$f_{2}-f_{1}-f_{3}=ma_{1}$\ccnum{⑥}，

设在 $t_{1}$ 时刻，B 与木板达到共同速度，其大小为 $v_{1}$，由运动学公式得

$v_{1}=v_{0}-a_{B}t_{1}$\ccnum{⑦}，$v_{1}=a_{1}t_{1}$\ccnum{⑧}，

联立\ccnum{①}\ccnum{②}\ccnum{③}\ccnum{④}\ccnum{⑤}\ccnum{⑥}\ccnum{⑦}\ccnum{⑧}式，代入已知数据得

$v_{1}=1\Ums$\ccnum{⑨}。

（2）在 $t_{1}$ 时间间隔内，B 相对于地面移动的距离为

$s_{B}=v_{0}t_{1}-\frac{1}{2}a_{B}t_{1}^{2}$\ccnum{⑩}，

设在 B 与木板达到共同速度 $v_{1}$ 后，木板的加速度大小为 $a_{2}$，对于 B 与木板组成的体系，由牛顿第二定律得

$f_{1}+f_{3}=(m_{B}+m)a_{2}$\ccnum{⑪}，

由\ccnum{①}\ccnum{②}\ccnum{④}\ccnum{⑤}式知，$a_{A}=a_{B}$；再由\ccnum{⑦}\ccnum{⑧}式知，B 与木板相反。由题意知，A 和 B 相遇时，A 与木板的速度相同，设其大小为 $v_{2}$。设 A 的速度大小从 $v_{1}$ 变到 $v_{2}$ 所用时间为 $t_{2}$，则由运动学公式，对木板，有

$v_{2}=v_{1}-a_{2}t_{2}$\ccnum{⑫}，

对 A，有

$v_{2}=-v_{1}+a_{A}t_{2}$\ccnum{⑬}，

在 $t_{2}$ 时间间隔内，B（以及木板）相对地面移动的距离为

$s_{1}=v_{1}t_{2}-\frac{1}{2}a_{2}t_{2}^{2}$\ccnum{⑭}，

在 $(t_{1}+t_{2})$ 时间间隔内，A 相对地面移动的距离为

$s_{A}=v_{0}(t_{1}+t_{2})-\frac{1}{2}a_{A}(t_{1}+t_{2})^{2}$\ccnum{⑮}，

A 和 B 相遇时，A 与木板的速度也恰好相同。因此 A 和 B 开始运动时，两者之间的距离为 $s_{0}=s_{A}+s_{1}+s_{B}$\ccnum{⑯}，

联立以上各式，并代入数据得

$s_{0}=1.9\Um$\ccnum{⑰}。

（也可以如图的速度$-$时间图线求解。）

\onepicture[w=7.0cm]{figs/25_an.png}
}

\item
%% number: 33
%% paperName: 2017年全国理综Ⅲ第 33 题 10 分
%% typeId: 6
%% type: 计算题
%% chapter: 气体性质
%% point: 玻意耳定律
%% method: 无
%% score: 10
%% degree: 450
%% duplicateId: 0
%% body:
\begin{enumerate}
	\item
	（5 分）如图，一定质量的理想气体从状态 a 出发，经过等容过程 ab 到达状态 b，再经过等温过程 bc 到达状态 c，最后经等压过程 ca 回到状态 a。下列说法正确的是\xzanswer{ABD}
	\fivechoices[answer=ABD]
	{在过程 ab 中气体的内能增加}
	{在过程 ca 中外界对气体做功}
	{在过程 ab 中气体对外界做功}
	{在过程 bc 中气体从外界吸收热量}
	{在过程 ca 中气体从外界吸收热量}
	\item
	（10 分）一种测量稀薄气体压强的仪器如图（a）所示，玻璃泡 M 的上端和下端分别连通两竖直玻璃细管 $K_{1}$ 和 $K_{2}$。$K_{1}$ 长为 $l$，顶端封闭，$K_{2}$ 上端与待测气体连通；M 下端经橡皮软管与充有水银的容器 R 连通。开始测量时，M 与 $K_{2}$ 相通；逐渐提升 R，直到 $K_{2}$ 中水银面与 $K_{1}$ 顶端等高，此时水银已进入 $K_{1}$，且 $K_{1}$ 中水银面比顶端低 $h$，如图（b）所示。设测量过程中温度、与 $K_{2}$ 相通的待测气体的压强均保持不变。已知 $K_{1}$ 和 $K_{2}$ 的内径均为 $d$，M 的容积为 $V_{0}$，水银的密度为 $\rho$，重力加速度大小为 $g$。求：
	\begin{enumerate}
		\item
		待测气体的压强；
		\item
		该仪器能够测量的最大压强。
	\end{enumerate}
\end{enumerate}
\onepicture[w=5.2cm, align=right]{figs/33a.png}
\twopicture[numA={（a）}, numB={（b）}, labelA=2017全国理综Ⅲ33b, labelB=2017全国理综Ⅲ33c, align=right, wA=4.8cm, wB=4.8cm]
{figs/33b.png}
{figs/33c.png}

%% answer: 
\jdanswer{
\begin{enumerate}
	\item
	ABD

	\item
	（i）$p_{x}=\frac{\pi\rho gh^{2}d^{2}}{4V_{0}+\pi d^{2}(l-h)}$；（ii）$p_{\max}=\frac{\pi\rho gl^{2}d^{2}}{4V_{0}}$

\end{enumerate}
}
%% memo:
\memoanswer{
（1）在过程 ab 中，体积不变，外界不对气体做功，气体也不对外界做功，压强增大，温度升高，内能增加，故 A 正确、C 错误；在过程 ca 中，气体的体积减少，外界对气体做功，故 B 正确；在过程 bc 中，温度不变，内能不变，体积增大，气体对外界做功，由热力学第一定律可知，气体要从外界吸收热量，故 D 正确；在过程 ca 中，压强不变，体积减小，温度降低，故内能减小，而外界对气体做功，气体要向外界放出热量，故 E 错误。

（2）（i）水银面上升至 M 的下端玻璃泡中气体恰好被封住，设此时被封闭的气体的体积为 $V$，压强等于待测气体的压强 $p$。提升 R，直到 $K_{2}$ 中水银面与 $K_{1}$ 顶端等高时，$K_{1}$ 中水银面比顶端低 $h$；设此时封闭气体的压强为 $p_{1}$，体积为 $V_{1}$，则

$V=V_{0}+\frac{1}{4}\pi d^{2}l$\ccnum{①}，$V_{1}=\frac{1}{4}\pi d^{2}h$\ccnum{②}，

由力学平衡条件得

$p_{1}=p+\rho gh$\ccnum{③}，

整个过程为等温过程，由玻意耳定律得

$pV=p_{1}V_{1}$\ccnum{④}，

联立\ccnum{①}\ccnum{③}\ccnum{④}式得

$p=\frac{\rho\pi gh^{2}d^{2}}{4V_{0}+\pi d^{2}(l-h)}$\ccnum{⑤}。

（ii）由题意知 $h\leq l$\ccnum{⑥}，

联立\ccnum{⑤}\ccnum{⑥}式有 $p\leq\frac{\pi\rho gl^{2}d^{2}}{4V_{0}}$\ccnum{⑦}，

该仪器能够测量的最大压强为 $p_{\max}=\frac{\pi\rho gl^{2}d^{2}}{4V_{0}}$\ccnum{⑧}。
}

\item
%% number: 34
%% paperName: 2017年全国理综Ⅲ第 34 题 10 分
%% typeId: 6
%% type: 计算题
%% chapter: 光学
%% point: 光的折射
%% method: 无
%% score: 10
%% degree: 550
%% duplicateId: 0
%% body:
\begin{enumerate}
	\item
	（5 分）如图，一列简谐横波沿 $x$ 轴正方向传播，实线为 $t=0$ 时的波形图，虚线为 $t=0.5\Us$ 时的波形图。已知该简谐波的周期大于 $0.5\Us$。关于该简谐波，下列说法正确的是\xzanswer{BCE}
	\fivechoices[answer=BCE]
	{波长为 $2\Um$}
	{波速为 $6\Ums$}
	{频率为 $1.5\UHz$}
	{$t=1\Us$ 时，$x=1\Um$ 处的质点处于波峰}
	{$t=2\Us$ 时，$x=2\Um$ 处的质点经过平衡位置}
	\item
	（10 分）如图，一半径为 $R$ 的玻璃半球，O 点是半球的球心，虚线 OO$'$ 表示光轴（过球心 O 与半球底面垂直的直线）。已知玻璃的折射率为 1.5。现有一束平行光垂直入射到半球的底面上，有些光线能从球面射出（不考虑被半球的内表面反射后的光线）。求：
	\begin{enumerate}
		\item
		从球面射出的光线对应的入射光线到光轴距离的最大值；
		\item
		距光轴 $\frac{R}{3}$ 的入射光线经球面折射后与光轴的交点到 O 点的距离。
	\end{enumerate}
\end{enumerate}
\twopicture[labelA=2017全国理综Ⅲ34a, labelB=2017全国理综Ⅲ34b, align=right, wA=8.5cm, wB=6.3cm]
{\includesvg{figs/34a}}
{figs/34b.png}

%% answer: 
\jdanswer{
\begin{enumerate}
	\item
	BCE

	\item
	（i）$\frac{2}{3}R$；（ii）$\frac{3(2\sqrt{2}+\sqrt{3})}{5}R\approx2.74R$

\end{enumerate}
}
%% memo:
\memoanswer{
（1）由波形图可知，波长为 $4\Um$，故 A 错误；简谐波沿 $x$ 轴正方向传播，实线为 $t=0$ 时的波形图，虚线为 $t=0.5\Us$ 时的波形图。又该简谐波的周期大于 $0.5\Us$，波传播的距离 $\Delta x=\frac{3}{4}\lambda$，$\frac{3}{4}T=0.5\Us$，故周期 $T=\frac{2}{3}\Us$，频率为 $1.5\UHz$，波速 $v=\lambda f=6\Ums$，故 B、C 正确；$t=1\Us=\frac{3}{2}T$，$t=0$ 时，$x=1\Um$ 处的质点处于波峰，$t=1\Us$ 时，该质点处于波谷，故 D 错误；$t=2\Us=3T$，是周期的整数倍，$t=0$ 时，$x=2\Um$ 处的质点位于平衡位置，$t=2\Us$ 时，该质点经过平衡位置，故 E 正确。

（2）（i）如图所示，从底面上 A 处射入的光线，在球面上发生折射时的入射角为 $i$，当 $i$ 等于全反射临界角 $i_{c}$ 时，对应入射光线到光轴的距离最大，设最大距离为 $l$。

$i=i_{c}$\ccnum{①}，

设 $n$ 是玻璃的折射率，由全反射临界角的定义得

$n\sin i_{c}=1$\ccnum{②}，

由几何关系得 $\sin i=\frac{l}{R}$\ccnum{③}，

联立\ccnum{①}\ccnum{②}\ccnum{③}式并利用题给条件得

$l=\frac{2}{3}R$\ccnum{④}。

\onepicture[w=6.0cm]{figs/34_an.png}

（ii）设与光轴相距 $\frac{R}{3}$ 的光线在球面 B 点发生折射时的入射角和折射角分别为 $i_{1}$ 和 $r_{1}$，由折射定律得

$n\sin i_{1}=\sin r_{1}$\ccnum{⑤}，

设折射光线与光轴的交点为 C，在 $\triangle OBC$ 中，由正弦定理得

$\frac{\sin\angle C}{R}=\frac{\sin(180^{\circ}-r_{1})}{OC}$\ccnum{⑥}，

由几何关系得 $\angle C=r_{1}-i_{1}$\ccnum{⑦}，$\sin i_{1}=\frac{1}{3}$\ccnum{⑧}，

联立\ccnum{⑤}\ccnum{⑥}\ccnum{⑦}\ccnum{⑧}式及题给条件得

$OC=\frac{3(2\sqrt{2}+\sqrt{3})}{5}R\approx2.74R$\ccnum{⑨}。
}

\end{enumerate}

\end{document}
